\chapteropen{07}

\begin{cele}
  \item Policzyć, ile dom realnie zjada w tygodniu, zanim cokolwiek w nim zmienisz.
  \item Ułożyć zakupy tak, żeby zajmowały jedną rundę w tygodniu, a nie codzienne wyjścia.
  \item Zaplanować gotowanie hurtowe na niedzielę i wiedzieć, kiedy jest tańsze niż gotowanie codzienne.
  \item Utrzymać czystość w progu dzięki zasadzie dziesięciu minut i jednej strefie dziennie.
  \item Podzielić obowiązki z domownikami i policzyć, czy podział się zgadza.
  \item Porównać koszt usługi z wartością godziny, którą naprawdę wykorzystasz.
\end{cele}

Dom jest najmniej widocznym konsumentem doby. Nie ma w nim jednego dużego
zadania, które wypada z kalendarza. Są dziesiątki drobnych czynności: lista
zakupów, której brakuje w połowie wieczoru, zmywarka, która czeka na
opróżnienie, pranie, które zapomniane w pralce pachnie już zupełnie
inaczej. Każda z nich zajmuje kilka minut, a łącznie potrafią zjeść dwie
godziny dziennie. W tym rozdziale nie będziemy ich ani ignorować, ani
udawać, że zniknął problem. Zamiast tego dom dostanie własną pozycję w
budżecie doby, a każda czynność domowa będzie musiała zmieścić się w
wyznaczonym limicie.

Warto od razu przyjąć jedną rzecz: nie zamierzamy wyprowadzić domu na
poziom, który wymagałby ciągłego wysiłku. Celem jest \emph{minimum}, czyli
stan, w którym dom jest wystarczająco czysty, jedzenie jest wystarczająco
zdrowe, a rzeczy wystarczająco łatwo znaleźć. Każdy ponadto kosztuje czas,
który w tym budżecie nie jest darmowy.

\section{Dom jako jedna pozycja w dobie}

W rozdziale 6 zbudowaliśmy szablony dnia. W każdym z nich pojawia się
blok \emph{Dom}, który celowo jest mały. W szablonie pracy zdalnej to
30 minut dziennie, w szablonie biura z dojazdem także 30 minut, a w
weekendowym szablonie 1,5 godziny na cały dzień. Liczby te nie są
wynikiem badań ani wzorcem dla wszystkich. Są decyzją projektową: budżet
musi zostawić miejsce na sen (\val{week.A.sleep} godzin), pracę i
pozostałe bloki, więc dom dostaje tyle, ile zostaje.

Trzydzieści minut dziennie to w dni robocze dwie i pół godziny w tygodniu.
Weekendowy blok 1,5 godziny dodaje kolejne półtorej, więc łącznie wychodzi
cztery godziny tygodniowo. To jest budżet, a nie rzeczywisty koszt:
prawdziwy koszt zmierzysz dopiero po dwóch tygodniach notowania, bo
pierwszy tydzień zawsze wygląda inaczej niż ten, który planowałeś.

Dlaczego notować? Bo większość ludzi przecenia to, co robi w domu, i
niedocenia tego, na co zużywa czas bez wyraźnego powodu. Zapis z telefonu,
choćby w postaci jednej linijki na dzień (zakupy 20 min, zmywanie 15 min,
pranie 5 min), wystarcza. Po dwóch tygodniach porównaj sumę z budżetem.
Jeśli wychodzi dwa razy więcej, nie zaczynaj od wymyślania nowych rozwiązań.
Zacznij od rozbicia wydatków czasowych na to, co konieczne, co
przyzwyczajenie i co zwyczajnie odkładasz.

\section{Zakupy: lista, dostawa i jedna runda w tygodniu}

Zakupy są największą pozycją w domu, która się da zredukować bez utraty
jakości życia. Zwykłe wyjście do sklepu w dni robocze to czas dojazdu,
chodzenia między regałami i kolejki przy kasie. Jeśli robisz to cztery razy
w tygodniu, płacisz za każde wyjście osobno. Jeśli robisz to raz, płacisz
raz.

Najprostszy układ opiera się na trzech elementach. Pierwszy to \emph{lista
stała}: kilkanaście do kilkudziesięciu pozycji, które kupujesz zawsze, z
ilością przy każdej. Lista nie musi być idealna, wystarczy, że pokrywa
około osiemdziesięciu procent koszyka. Drugi to \emph{dostawa}: zamówienie
z listy przychodzi pod drzwi, zamiast zajmować wyjście. Trzeci to
\emph{jedna większa runda tygodniowo}, na którą przeznaczasz konkretny
czas i do której wrzucasz to, czego nie ma na liście albo co kupuje się
tylko osobiście (świeże zioła, pieczywo, rzeczy, które trzeba obejrzeć).

\begin{przyklad}[Lista z dostawą w praktyce]
Załóżmy listę 25 pozycji, którą uzupełniasz raz w tygodniu w 20 minut,
w sobotę rano, przy kawie. Zamówienie składasz od razu, z wybranym oknem
dostawy. Czas domowy na zakupy to wtedy 20 minut tygodniowo zamiast
sześciu wyjść po 25 minut, czyli około dwóch i pół godziny. Różnica
wynosi więc ponad dwie godziny tygodniowo, zanim policzysz dojazd.
Ile dokładnie kosztuje dostawa w twoim sklepie, zależy od progu zamówienia
i od tego, czy masz subskrypcję. \verify{typowe opłaty za dostawę zakupów
spożywczych i progi kwotowe zmieniają się w czasie i między sklepami;
sprawdź cennik konkretnego sklepu przed decyzją}.
\end{przyklad}

Dostawa ma jedną wadę, której nie wolno zlekceważyć: łatwo do koszyka
wpadają rzeczy, których na liście nie było. Trzeba się pilnować
\emph{progu}: jeśli dostawa wymaga zamówienia za określoną kwotę, a lista
jest krótka, dopełnianie koszyka do progu jest zakupem impulsywnym. Lepiej
dopłacić drobną opłatę, niż dopełniać koszyk czymś, czego nie planowałeś.

Automatyzacja działa w tym samym miejscu. Subskrypcja na powtarzalne
produkty (kawa, makaron, ryż, środki czystości) ogranicza decyzje do
zatwierdzenia terminu. Przypomnienie na telefonie w czwartek wieczorem,
\glqq sprawdź listę przed sobotą\grqq{}, to koszt kilku sekund. Oba
rozwiązania działają tylko wtedy, gdy lista jest stała. Subskrypcja na
produkt, którego już nie jesz, to dokładnie ten rodzaj automatu, który
pracuje przeciwko budżetowi.

\section{Gotowanie hurtowe w niedzielę}

Gotowanie codzienne ma dwa koszty: czas na przygotowanie każdego posiłku
i czas na sprzątanie po nim. Gotowanie hurtowe przenosi oba koszty na
jedną sesję, w której robisz kilka dań naraz i dzielisz je na porcje.
Sesja w niedzielę jest wygodna, bo w tygodniu nie masz na nią miejsca.

Działa to pod jednym warunkiem: dania muszą znieść przechowywanie i ponowne
podgrzanie. Zupy, gulasze, sosy, kasze, strączki, pieczone warzywa i
mięso w sosie zwykle się sprawdzają. Sałaty, smażone panierki i wszystko,
co ma być świeże, nie. Ta granica decyduje o tym, co wpisujesz do planu.

\begin{przyklad}[Niedzielna sesja na dziesięć porcji]
Plan: jedna zupa, jeden gulasz z kaszą, jedna porcja pieczonych warzyw na
dwa dni i jedna potrawa z roślin strączkowych. Składniki kupujesz na
liście z soboty. Czas aktywny: około godziny i dziesięciu minut, bo część
pracy wykonuje piekarnik i garnek. Sprzątanie po sesji: kwadrans, jeśli
naczynia myjesz od razu. Razem około półtorej godziny, czyli dokładnie
tyle, ile blok \emph{Dom} w weekendowym szablonie.

Te same dziesięć porcji zrobione codziennie zajęłyby kilkakrotnie więcej
czasu aktywnego, bo każde gotowanie wymaga dojścia do garnka, przypraw i
mycia. Różnica nie jest więc wyłącznie w czasie, lecz w liczbie
\emph{startów}, a start jest najdroższą częścią gotowania.
\end{przyklad}

Pamiętaj o przechowywaniu. Porcje, które mają czekać dłużej niż kilka dni,
powinny trafić do zamrażarki, a te, które zjesz w ciągu dwóch dni, do
lodówki w szczelnych pojemnikach. \verify{zalecany czas przechowywania
gotowanych potraw w lodówce i zamrażarce zależy od rodzaju potrawy i
podawany jest w wytycznych dotyczących bezpieczeństwa żywności; sprawdź
aktualne wytyczne, zanim ustalisz własne terminy}.

Jeśli sesja trwa dłużej niż zakładasz, nie wydłużaj jej do wieczora.
Lepiej zrobić mniej dań, a resztę przenieść na następną niedzielę. Hurtowe
gotowanie ma sens tylko jako stała, przewidywalna pozycja; sesja, która
co tydzień wypada poza budżet, jest kolejnym źródłem presji.

\section{Sprzątanie w progu i pranie}

Sprzątanie kojarzy się z jednym dużym wysiłkiem: sobota, odkurzacz, mycie
wszystkiego naraz. W budżecie doby ten model nie działa, bo jedna sobota
pochłania godziny, a zanim przyjdzie następna, dom znów wymaga uwagi. Lepiej
rozłożyć pracę tak, żeby każdego dnia wykonać małą część.

Zasada dziesięciu minut mówi: codziennie masz dziesięć minut na jedną
konkretną czynność, a po ich upływie kończysz, nawet jeśli nie skończyłeś.
Nie jest to pozwolenie na bałagan. Jest to ograniczenie, które wymusza
wybór. Dziesięć minut wystarcza na zmycie naczyń z dnia, wytarcie blatów,
odłożenie rzeczy z podłogi albo wyczyszczenie jednego obszaru łazienki.

Drugi element to \emph{jedna strefa dziennie}. Dom dzielisz na kilka
stref, na przykład kuchnię, łazienkę, przedpokój, salon, sypialnię, pokój
dzienny i przestrzeń wspólną. Każdego dnia wybierasz jedną strefę i
sprzątasz ją w dziesięć minut zamiast czterdziestu. Przy siedmiu strefach
każda jest sprzątana raz w tygodniu, a łączny koszt to 70 minut, a nie cała
sobota.

\begin{przyklad}[Siedem stref w tygodniu]
Strefy: kuchnia, łazienka, przedpokół, salon, sypialnia, biurko, przestrzeń
wspólna. Poniedziałek sprzątasz kuchnię, wtorek łazienkę i tak dalej.
Jeśli w danym dniu wypadnie ci dziesięć minut dodatkowego czasu, nie
przechodzisz do następnej strefy, tylko zostawiasz ją na później. Dziesięć
minut na strefę razy siedem stref to 70 minut tygodniowo, czyli mniej niż
połowa jednego wieczoru.
\end{przyklad}

Pranie wymaga innego podejścia, bo większość jego czasu to czas bierny.
Maszyna pracuje, a ty nie musisz przy niej stać. Wystarczy załadowanie,
uruchomienie i rozwieszenie. Załadowanie zajmuje kilka minut, rozwieszenie
kolejne kilka. Kluczowe jest, żeby pranie nie zalegało: pranie, które
czeka trzy dni, zajmuje dwa razy więcej czasu, bo trzeba je przeprasować
albo przemyć ponownie.

Ustaw rytm: jedno pranie dziennie albo co drugi dzień, zawsze w tym samym
czasie. Pralka działa, kiedy ty jesz kolację, a rano rozwieszasz. Program
trwa zwykle godzinę lub dłużej, \verify{czas trwania programów pralki i
suszenia zależy od modelu urządzenia i wybranego programu; sprawdź
instrukcję swojego sprzętu}. Nie zostawiaj włączonej suszarki bez nadzoru,
kiedy idziesz spać, i nie traktuj pralki jako urządzenia, które bezpiecznie
pracuje w nocy bez kontroli; to kwestia bezpieczeństwa, nie wygody.

\section{Podział obowiązków z domownikami}

Jeśli mieszkasz z kimś, budżet domowy jest wspólny, nawet gdy konta są
osobne. Każda godzina, którą jedna osoba przeznacza na dom, jest godziną,
której druga nie przeznacza na swój sen albo pracę. Nierównowaga zwykle nie
wynika ze złej woli. Wynika z tego, że nikt nie policzył, kto co robi, a
ten, kto robi więcej, przestaje o tym mówić.

Rozmowa powinna zacząć się od liczb, nie od oceny. Zapisz wszystkie
obowiązki, które pojawiają się w domu, z częstotliwością i szacowanym
czasem. Potem przeczytajcie listę razem i wspólnie zdecydujcie, co jest
stałe, co jest jednorazowe, a co można przerzucić na usługę albo
automat. Dopiero po tym przydzielcie zadania, najlepiej tak, żeby każdy
miał jedną stałą pozycję i jedną elastyczną.

\begin{tabular}{@{}lllr@{}}
\toprule
\textbf{Zadanie} & \textbf{Kto} & \textbf{Jak często} & \textbf{Minuty/tydz.}\\
\midrule
Lista zakupów i zamówienie & osoba A & raz w tygodniu & 20\\
Gotowanie hurtowe & osoba B & raz w tygodniu & 90\\
Strefa dzienna & na zmianę & codziennie & 70\\
Pranie (załadunek i rozwieszenie) & osoba B & co drugi dzień & 30\\
Wynoszenie śmieci i segregacja & osoba A & co tydzień & 15\\
\bottomrule
\end{tabular}

Tabela nie jest sprawiedliwa z definicji. Jest tylko jawna. Jeśli jedna
osoba ma 180 minut tygodniowo, a druga 40, nie znaczy to jeszcze, że
pierwsza jest przeciążona. Może pracuje mniej w domu zawodowo, może ma
dłuższe dojazdy. Chodzi o to, żeby różnicę dało się obejrzeć i omówić, a
nie żeby wynik był równy co do minuty.

Podział trzeba też przeglądać. Co miesiąc sprawdźcie, czy tabela nadal
odpowiada rzeczywistości, bo zmiana pracy, dziecko albo przeprowadzka
przesuwają ciężar bez niczyjej decyzji.

\section{Pieniądze zamiast czasu}

Każdą czynność domową można kupić. Sprzątanie firmowe, dostawa zakupów,
gotowe posiłki, pralnia, usługa montażu mebli: wszystko to zamienia
pieniądze na czas. Decyzja, czy warto, nie jest moralna ani ideologiczna.
Jest rachunkiem stawki godzinowej.

Krok pierwszy to stawka \emph{z czasu, który realnie wykorzystasz}. Nie
bierz pod uwagę hipotetycznej stawki z pensji, jeśli za każdą zaoszczędzoną
godzinę i tak zrobisz coś, co nie ma wartości. Zapytaj, czym zajęłabyś lub
zająłbyś tę godzinę: snem, odpoczynkiem, pracą płatną albo czymś, co
rzeczywiście cię cieszy. Jeśli odpowiedź brzmi \glqq przewinę telefon\grqq{},
wartość godziny jest niska i nie kupujesz jej po wysokiej cenie.

Krok drugi to koszt usługi podzielony przez zaoszczędzony czas.

\begin{przyklad}[Stawka godzinowa usługi]
Przyjmijmy, że godzina twojego czasu ma dla ciebie wartość 100 zł. Jest to
założenie do policzenia, nie wynik z pensji; w swoim przypadku wstaw liczbę,
w którą naprawdę wierzysz. Usługa sprzątania trwa trzy godziny i kosztuje
240 zł. \verify{typowe ceny sprzątania domu w twojej okolicy różnią się
w zależności od metrażu, zakresu i firmy; wstaw cenę z własnej oferty}.

Koszt jednej zaoszczędzonej godziny: 240 zł podzielone przez 3 godziny,
czyli 80 zł. Wartość godziny to 100 zł. Usługa wychodzi taniej niż twój
czas, więc na papierze się opłaca.

Ten sam rachunek dla dostawy zakupów: opłata 20 zł za zamówienie, które oszczędza ponad dwie godziny, daje około 9 zł za godzinę oszczędzonego czasu. Jeśli twoja godzina ma wartość 100 zł, to jest tanio. Ale jeśli wartość godziny spada do 8 zł, bo wieczór i tak mija na przewijaniu, ta sama opłata przestaje mieć sens, mimo że sama w sobie jest niska.
\end{przyklad}

Dwie zasady wynikają z tego rachunku. Po pierwsze, usługa jest opłacalna,
gdy jej koszt za godzinę jest niższy od wartości godziny, którą naprawdę
wykorzystasz. Po drugie, opłaca się przede wszystkim to, co zdejmuje z
harmonogramu czynności stałe i nieprzyjemne, których nie da się skrócić.
Mycie okien co miesiąc to dobry kandydat, a zmywanie po kolacji już nie.

Rachunek ma jedną pułapkę. Godziny zaoszczędzone w kawałkach po dziesięć
minut nie układają się w godzinę, którą można sprzedać albo zaplanować.
Usługa, która oszczędza czterdzieści minut tygodniowo, rzadko zmienia coś
konkretnego, chyba że zdecydujesz, co zrobisz z tym czasem. Zanim
zapłacisz, zapisz, co dokładnie zajmie ten czas.

\section{Pułapka perfekcjonizmu}

Najczęstszy powód, dla którego plan domowy się rozpada, nie jest brak czasu.
Jest nim przekonanie, że jeśli sprzątanie nie jest dokładne, to nie ma
sensu go zaczynać. Dziesięć minut wydaje się wtedy zbyt małe, żeby coś
zmienić, więc nie robi się nic. Potem dochodzi poczucie winy, a z nim
kolejna godzina wypadająca z planu.

Budżet doby jest napisany przeciw perfekcjonizmowi. Trzydzieści minut
dziennie nie pozwala na czystą kuchnię, idealną łazienkę i wypolerowane
podłogi. Pozwala na dom, który jest wystarczająco czysty, żeby w nim żyć,
gotować i wracać po pracy bez zniechęcenia. Pełne sprzątanie zajęłoby
kilkakrotnie więcej czasu, a w zamian dałoby jedynie różnicę, której
większość gości nie zauważy.

\begin{falle}[Dom musi być czysty, zanim pójdę spać]
Ten warunek zawsze kończy się późnym wieczorem, zmęczeniem i przesunięciem
snu, czyli dokładnie tym, co miało budżet chronić. Poprawka: wieczorem
wykonujesz jedną czynność z listy na dziesięć minut, a reszta czeka do
następnego dnia. Dom nie musi być w stanie idealnym w chwili, w której
kładziesz się spać. Musi być w stanie, z którego rano da się wyjść bez
szukania klucza.
\end{falle}

Perfekcjonizm ma też drugą formę: odkładanie decyzji, bo żadna z opcji nie
jest idealna. Zamiast wybrać lepszą listę zakupów, przez tydzień
przerzucasz trzy aplikacje. Zamiast kupić dostawę, bo jest wystarczająca,
czekasz na promocję, której nie ma. Zasada: jeśli rozwiązanie spełnia
zapisane kryteria, wdrażasz je i przeglądasz za miesiąc. Nie zakładasz, że
trzeba je dopracować przed startem.

Warto też zapytać, co w ogóle jest celem. Dom w minimum nie jest
projektem estetycznym. Jest warunkiem, żeby reszta planu działała: żeby
sen nie zaczynał się od zmywania, żeby kolacja nie wymagała godziny
przygotowań, żeby zakupy nie wypadały w poniedziałek o dwudziestej.
Jeśli dom spełnia tę funkcję, reszta jest decyzją, nie obowiązkiem.

\begin{zasada}[Minimum wystarczy, dopóki działa]
Dom ma spełniać trzy warunki: da się w nim spać, da się w nim gotować i da
się znaleźć potrzebne rzeczy. Wszystko, co wykracza poza to, jest wyborem
i musi zmieścić się w budżecie, nie kosztem snu.
\end{zasada}

\begin{cwiczenia}
  \exitem{Ile godzin tygodniowo wynosi budżet bloku \emph{Dom} w szablonie pracy zdalnej i w szablonie weekendowym, jeśli oba są stosowane przez cały tydzień?}{W dni robocze szablon A ma 0,5 godziny Domu, czyli pięć dni daje 2,5 godziny. Szablon weekendowy ma 1,5 godziny, więc łącznie wychodzi 2,5 + 1,5 = 4 godziny tygodniowo.}
  \exitem{Lista zakupów ma 25 pozycji i zajmuje 20 minut raz w tygodniu. Sześć wyjść do sklepu po 25 minut każde zajmuje ile? Ile minut oszczędza zamówienie?}{Sześć wyjść po 25 minut to 150 minut, czyli dwie i pół godziny. Zamówienie zajmuje 20 minut, więc oszczędność wynosi 150 − 20 = 130 minut, czyli ponad dwie godziny tygodniowo.}
  \exitem{Gotowanie hurtowe daje dziesięć porcji w sesji trwającej około półtorej godziny. Ile porcji wypada na godzinę sesji i dlaczego to porównanie nie jest pełne?}{Dziesięć porcji podzielone przez półtorej godziny daje około 6,7 porcji na godzinę. Porównanie jest niepełne, bo nie uwzględnia czasu przechowywania ani kosztu składników; liczy tylko czas sesji.}
  \exitem{Zasada dziesięciu minut i siedem stref. Ile minut tygodniowo zajmuje pełny cykl, jeśli każdego dnia sprzątasz jedną strefę? Jaki wniosek wynika z tego dla soboty?}{Siedem stref po 10 minut to 70 minut tygodniowo. Wniosek: sobota nie musi być dniem sprzątania całego domu, bo cały cykl mieści się w jednym krótkim wieczorze tygodnia.}
  \exitem{Usługa kosztuje 240 zł i trwa trzy godziny. Wartość godziny w twoim rachunku to 100 zł. Czy usługa się opłaca i od czego zależy odpowiedź?}{Koszt za godzinę to 240 : 3 = 80 zł, poniżej wartości 100 zł, więc na papierze się opłaca. Odpowiedź zależy od tego, czy zaoszczędzony czas zostanie wykorzystany na coś, co ma wartość; jeśli trafi do przewijania, rachunek nie działa.}
  \exitem{Podział obowiązków: osoba A ma 180 minut tygodniowo, osoba B 40. Czy to dowód nierówności? Co trzeba sprawdzić, zanim wyciągniesz wniosek?}{Sama różnica nie dowodzi nierówności. Trzeba sprawdzić, ile czasu każda osoba ma na obowiązki zawodowe, dojazdy i inne zobowiązania oraz czy tabela odpowiada rzeczywistości.}
  \exitem{Dlaczego pranie zalegające trzy dni kosztuje więcej czasu niż pranie co drugi dzień?}{Trzeba je ponownie przemyć, przeprasować albo wyprać z innym programem, a każdy dodatkowy cykl to kolejne załadowanie i rozwieszenie.}
\end{cwiczenia}

\zadanie{Przez dwa tygodnie zapisuj czas, który realnie zjada dom, w jednej linijce na dzień
(zakupy, gotowanie, sprzątanie, pranie). Po drugim tygodniu zsumuj każdą kategorię i porównaj
z budżetem czterech godzin tygodniowo. Miara sukcesu: masz w jednym zdaniu wskazaną największą
pozycję, której różnica względem budżetu jest najbardziej widoczna, i jedną zmianę, którą wprowadzisz
w kolejnym tygodniu, z przewidywanym czasem oszczędności.}
